%%#! latex

%%%%%%%%%%             Preamble                   %%%%%%%%%%

\documentclass[10pt]{jarticle}
\usepackage{amsfonts,amssymb,amsmath,amscd,amsthm,latexsym}
\usepackage{graphicx}
\usepackage{fancybox}
\usepackage{okumacro}
%\usepackage{hyperref}
%\usepackage{showkeys}
%\usepackage{fancybox,ascmac}
%\usepackage{fancyhdr}
\pagestyle{empty}

%%%%%%%%%%             Layout                     %%%%%%%%%%

\setlength{\topmargin}{-1cm}
\setlength{\textheight}{25cm}
\setlength{\textwidth}{16cm}
%\setlength{\oddsidemargin}{0cm}
%\setlength{\evensidemargin}{0cm}
%\addtolength{\textwidth}{.815in}
\addtolength{\oddsidemargin}{-2cm}
\addtolength{\evensidemargin}{-2cm}

%%%%%%%%%%           Definition, Theorem, etc.     %%%%%%%%%%

\newtheorem{thm}{定理}
\newtheorem{lem}[thm]{補題}
\newtheorem{prop}[thm]{命題}
\newtheorem{claim}[thm]{主張}
\newtheorem{cor}[thm]{系}
\newtheorem{prob}[thm]{問題}
\newtheorem{conj}[thm]{予想}
\theoremstyle{definition}
\newtheorem{dfn}[thm]{定義} 
\newtheorem{ex}[thm]{例}
\newtheorem{reidai}[thm]{例題}
\newtheorem{rmk}[thm]{注意}

%\numberwithin{thm}{section}
%\numberwithin{equation}{section}

%%%%%%%%%%             Macro (by Nasu)           %%%%%%%%%%

\newcommand{\Proof}{\noindent {\bf 証明)}\ \ }
\newcommand{\Hom}{\operatorname{Hom}}
\newcommand{\Ext}{\operatorname{Ext}}
\newcommand{\Spec}{\operatorname{Spec}}
\newcommand{\Proj}{\operatorname{Proj}}
\newcommand{\ob}{\operatorname{ob}}
\newcommand{\Hilb}{\operatorname{Hilb}}
\newcommand{\red}{\operatorname{red}}
\newcommand{\im}{\operatorname{im}}
\newcommand{\Bs}{\operatorname{Bs}}
\newcommand{\Jac}{\operatorname{Jac}}
\newcommand{\Pic}{\operatorname{Pic}}
\newcommand{\coker}{\operatorname{coker}}
\newcommand{\Gr}{\operatorname{Gr}}
\newcommand{\Bl}{\operatorname{Bl}}
%\newcommand{\algeeq}{\stackrel{alg.}{\; \sim \;}}
\newcommand{\Rat}{\operatorname{Rat}}
\newcommand{\Exc}{\operatorname{Exc}}
\newcommand{\car}{\operatorname{char}}
\newcommand{\Aut}{\operatorname{Aut}}
\newcommand{\End}{\operatorname{End}}
\newcommand{\sat}{\operatorname{sat}}
\newcommand{\pd}{\operatorname{dp}}
\newcommand{\rank}{\operatorname{rank}}
\newcommand{\Ann}{\operatorname{Ann}}

\renewcommand{\theenumi}{\arabic{enumi}}
\renewcommand{\theenumii}{\alph{enumii}}
\renewcommand{\theenumiii}{\roman{enumiii}}
\renewcommand{\labelenumi}{(\theenumi)}
\renewcommand{\labelenumii}{(\theenumii)}
\renewcommand{\labelenumiii}{[\theenumiii]}
%\renewcommand{\baselinestretch}{1.2}
%\renewcommand{\abstractname}{\sc Abstract}

%%%%%%%%%%           Commutative diagram          %%%%%%%%%%

% A litte smaller diagram than AMS CD envrionment

\newcommand{\mapright}[1]{%
\smash{\mathop{%
\hbox to 1cm{\rightarrowfill}}\limits^{#1}}}
\newcommand{\mapleft}[1]{%
\smash{\mathop{%
\hbox to 1cm{\leftarrowfill}}\limits^{#1}}}
\newcommand{\mapdown}[1]{\Big\downarrow
\llap{$\vcenter{\hbox{$\scriptstyle#1\,$}}$ }}
\newcommand{\mapup}[1]{\Big\uparrow
\rlap{$\vcenter{\hbox{$\scriptstyle#1$}}$ }}

%%%%%%%%%%              Title                    %%%%%%%%%%

\title{}
\author{}
\date{}

%%%%%%%%%%            Text Start                 %%%%%%%%%%

\begin{document}

%\maketitle

\begin{center}
 {\bf \Huge 東海大学理学部数学・情報数理談話会}
\end{center}

\bigskip

{\Large
以下の要領において談話会を開催致します. 多数の方の御来聴を
お待ち致しております. 

\bigskip

\noindent
日時: 2026年6月26日(金) 16:00--17:00\\
\smallskip
会場: 東海大学湘南校舎18号館{\bf 7階情報数理ゼミ室1 (18-701)}\\

\medskip

\noindent
講演者: 
安保広達氏(Department of Mathematics and Statistical Science, University of Idaho)\\
\smallskip
タイトル: 
On the classification of surfaces of non-general type in projective four-space\\

\noindent
アブストラクト:
Hartshorne and Lichtenbaum conjectured that the degrees
of smooth rational surfaces in projective four-space are bounded. In 1989,
Ellingsrud and Peskine confirmed the conjecture by proving a stronger
result: the family of smooth surfaces of non-general type in projective
four-space is bounded.

The classification of smooth surfaces in projective four-space has been
completed up to degree 10. Significant progress has been made toward
classifying smooth, non-general-type surfaces of degree 11. However, the
classification problem for higher-degree surfaces remains less well under-
stood.

This talk aims to review the current status of the classification of these
surfaces and to discuss previously unknown examples with degree 12, sec-
tional genus 13, and Euler characteristic 2.

This is joint work (still ongoing) with Kristian Ranestad and Frank-Olaf
Schreyer.

\vskip 2cm

\begin{flushright}
  世話人: 那須 弘和\\
  nasu@tokai.ac.jp
\end{flushright}
}
\end{document}

